% !TEX root = creche.tex
\documentclass[creche.tex]{subfiles}
\begin{document}
\chapter{Measures???}
\label{measures}

\def\medrel#1{\parbox[t]{5ex}{$\displaystyle\hfil #1$}}
\def\ceq#1#2#3{\parbox[t]{18ex}{$\displaystyle #1$}\medrel{#2}{$\displaystyle #3$}}


\def\Msr{{\rm Msr\kern.0ex}}


Scratch paper

\section{Keisler measures}

Fix a first-order structure $M$ and some tuple of variables $x$.
Let $\Delta\subseteq L_x(M)$ be a set of formulas closed under Boolean connectives.
A \emph{Keisler measure\/} on $\Delta$ is a function $\mu:\Delta\to\RR$ such that
\begin{itemize}
\item[1.] $\mu(\top)=1$;
\item[2.] $\mu(\phi\vee\psi)=\mu(\phi)+\mu(\psi)$ for every $\phi,\psi\in \Delta$ such that $M\models\neg\E x\,(\phi\wedge\psi)$.
\end{itemize}
When $\Delta=L_x(M)$ we say that $\mu$ is a \emph{global Keisler measure}.

% 
% To a global Keisler measure $\mu$ we associate a structure $M_\mu$ in continuous logic.
% Namely, $M$ itself is identified with a structure in continuous logic in the canonical way and $M_\mu$ is the expansion of $M$ that has a $[0,1]$-valued relation for every $\phi(x\,;z)\in L$
% 
% \ceq{\hfill \phi(\mu\,;\cdot)\ :\ M^{|z|}}{\mapsto}{[0,1]}
% 
% \ceq{\hfill b\ \ }{\mapsto}{\mu(\phi(x\,;b)\big)}
% 
% We write $L^{\rm c}$ for the expanded signature.

Let $M$ and $N$ be two structures with a common subset $A$.
Let $\mu$ and $\nu$ be two global Keisler measures on $M$ and $N$ respectively.
We say that $M_\mu\equiv_A N_\nu$ 

$\mu\big(\phi(x\,;b)\big)=\nu\big(\phi(x\,;b)\big)$ for every $b\in A^{|z|}$.


and let $A\subseteq M\cap N$.
 be a common fragment.


\section{Borel measures}

A Keisler measure on $\Delta\subseteq L_x(\U)$ induces a finite probability measure on the subsets of $\U^{|x|}$ that are $\Delta$-definable.
This in turn can be extended to a Borel measure on $\U^{|x|}$.
When $\Delta$ ha small cardinality, it is a regular Borel measure.
On the other hand, global Keisler measures may not always induce a regular measure.

In this section we discuss the extension of $\mu$ to a Borel measure.  

In this section by \emph{definable\/} we intend definable by a formula in $\Delta$.
A \emph{type-definable set\/} is the intersection of a small collection of definable sets.
 A \emph{type-codefinable set\/} is the complement of a type-definable set.

When $\Delta$ is small, type-definable and type-codefinable sets are the closed, respectively open, sets of the compact topology generated by the definable sets.

To begin with, we extend $\mu$ by defining it on any type-definable set $\F$ and type-codefinable set $\O$ 

\ceq{\hfill\mu(\F)}{=}{\ \,\inf\bigg\{\mu(\D)\ :\ \F\subseteq\D,\ \D\ \mathrm{definable}\bigg\}.}

\ceq{\hfill\mu(\O)}{=}{\sup\bigg\{\mu(\D)\ :\ \D\subseteq\O,\ \D\ \mathrm{definable}\bigg\};}

The two definitions agree on sets that are both type-definable and type-codefinable as these are, by compactness, definable.

\begin{proposition}\label{prop_measure_on_open_closed}
Additivity (i.e.\@ axiom \ssf{2} above) holds both among type-definable and among type-codefinable sets.
Moreover $\mu(\O\sm\F)=\mu(\O)-\mu(\F)$ for every type-definable set $\F$ contained in a type-codefinable set $\O$.
\end{proposition}
\begin{proof}
By compactness, any two disjoint type-definable  sets are separated by a definable set, so $\mu$ is additive on type-definable  sets.
As $\mu(\O)=1-\mu(\neg\O)$, it is also additive among type-codefinable set.

Now, let $\F\subseteq\O$ be as in the proposition.
Fix $\epsilon>0$ and let $\D\subseteq\O\sm\F$ be a definable set such that $\mu(\O\sm\F)-\epsilon\le\mu(\D)$.
By compactness, there is a definable set $\C$ containing $\F$ and disjoint from $\D$.
Then

\ceq{\hfill\mu(\O\sm\F)-\epsilon}{\le}{\mu(\D)}
\medrel{\le}
$\mu(\O\sm\C)$
\medrel{=}
$\mu(\O)-\mu(\C)$
\medrel{\le}
$\mu(\O)-\mu(\F)$ 

As $\epsilon$ is arbitrary, we obtain $\mu(\O\sm\F)\le\mu(\O)-\mu(\F)$.
For the converse inequality let $\D\supseteq\F$ such that  $\mu(\D)-\mu(\F)\le\epsilon$.
Then

\ceq{\hfill\mu(\O\sm\F)}{\ge}{\mu(\O\sm\D)}
\medrel{=}
$\mu(\O)-\mu(\D)$
\medrel{\ge}
$\mu(\O)-\mu(\F)-\epsilon$.

As $\epsilon$ is arbitrary, the converse inequality follows.\end{proof}

Now, for any subset $\X\subseteq\U^{|x|}$ we define the \emph{outer\/} and the \emph{inner\/} measure of $\X$ to be respectively

\ceq{\hfill\mu^*(\X)}{=}{\,\ \inf\bigg\{\mu(\O)\ :\ \X\subseteq\O,\ \O\ \textrm{type-codefinable}\bigg\};}

\ceq{\hfill\mu_*(\X)}{=}{\sup\bigg\{\mu(\F)\ :\ \F\subseteq\X,\ \F\ \textrm{type-definable}\bigg\}\medrel{=}1-\mu^*(\X).}

Clearly $\mu^*$ and $\mu_*$ coincide with $\mu$ on the definable sets.
We prove that they coincide also also on type-definable and type-codefinable sets.

\begin{proposition}
If $\X$ is type-definable or type-codefinable then $\mu^*(\X)=\mu_*(\X)=\mu(\X)$.
\end{proposition}
\begin{proof}
Let $\F$ be type-definable.
Then $\mu_*(\F)=\mu(\F)$ is obvious.
By compactness, for every type-codefinable set $\F\subseteq\O$, there is a definable set $\F\subseteq\D\subseteq\O$.
Then $\mu^*(\F)=\mu(\F)$ follows.
As $\mu_*(\X)=1-\mu^*(\neg\X)$, the same identities hold for type-codefinable sets.
\end{proof}

From the additivity of $\mu$ we easily infer that for every pair of sets $\X,\Y\subseteq\U^{|x|}$ such that $\X\cap\Y=\0$

\ssf{\#}\hfil$\mu_*(\X)+\mu_*(\Y)\medrel{\le}\mu_*(\X\cup\Y)\medrel{\le}\mu^*(\X\cup\Y)\medrel{\le}\mu^*(\X)+\mu^*(\Y)$

In words we say that the outer measure is \emph{subadditive\/} and that the inner measure is \emph{superadditive}.

From subadditivity it immediately follows that $\mu^*(\X\simdiff\Y)$ defines a pseudometric on $\P\big(\U^{|x|}\big)$.
Subadditivity also implies that $\mu^*(\X)-\mu^*(\Y)\le\mu^*(\X\simdiff\Y)$, therefore the function $\X\mapsto\mu^*(\X)$ is continuous, w.r.t.\@ this pseudometric.
It is also easy to verify that the functions $\X\mapsto\mu^*(\neg\X)$ and $\<\X,\Y\>\mapsto\mu^*(\X\cup\Y)$ are continuous.

As $\mu_*(\X)=1-\mu^*(\neg\X)$, the inner measure is also continuous.

We say that $\X$ is \emph{measurable\/} if it is in the closure of the definable sets in the topology induced by this pseudometric above. When $\X$ is measurable, we write \emph{$\mu(\X)$\/} for  $\mu^*(\X)$.

\begin{proposition}\label{prop_measurable_equivalents}
For every $\X$ the following are equivalent
\begin{itemize}
\item[1.] $\mu^*(\X)=\mu_*(\X)$;
\item[2.] $\X$ is measurable.
\end{itemize}
\end{proposition}
\begin{proof}
\ssf{1}$\IMP$\ssf{2}\quad First note that type-definable sets are measurable.
Now, fix $\epsilon>0$ and let $\F\subseteq\X$ be a type-definable set such that $\mu_*(\X)-\mu(\F)\le\epsilon$.

Using Proposition~\ref{prop_measure_on_open_closed} it is not difficult to show that $\mu^*(\X\sm\F)\le\mu^*(\X)-\mu(\F)$.
So, from \ssf{1} we obtain  $\mu^*(\X\simdiff\F)<\epsilon$.
As $\epsilon$ is arbitrary \ssf{2} follows.

\ssf{2}$\IMP$\ssf{1}\quad It suffices to prove that $\mu^*(\X)=1-\mu^*(\neg\X)$.
This identity holds for all definable set.
By continuity, it also holds for all measurabe sets.
\end{proof}

Let $\<r_i : i<\lambda\>$ be a sequence of non negative real numbers. 
We define

\ceq{\hfill\sum_{i<\lambda} r_i}{=}{\sup\Big\{\sum_{i\in I} r_i\ :\ I\subseteq\lambda\  \textrm{finite}\Big\}.}

Clearly, the sum is finite only $r_i=0$ for all but countably may $i$. 


\begin{proposition}
Let $\lambda$ be a small cardinal.
The union of an arbitrary collection $\<\X_i : i<\lambda\>$ of measurable sets is measurable. 
Moreover, if the $\X_i$ are pairwise disjoint

\ceq{\hfill\mu\bigg(\bigcup_{i<\lambda}\X_i\bigg)}{=}{\sum_{i<\lambda}\mu(\X_i)}
\end{proposition}

\begin{proof}
It is immediate that measurable sets are cosed under complementation.
By Proposition~\ref{prop_measurable_equivalents} and the inequalities \ssf{\#}, they are also closed under union.
Hence they form a boolean algebra.
We prove that measurable sets are closed under infinite union simultaneously with the second claim of the proposition.
We reason by induction on $\lambda$.
Clearly, only limit stages are left to be proved.
By induction hypothesis we can assume that the $\X_i$ are pairwise disjoint.
By Proposition~\ref{prop_measurable_equivalents} it suffices to verify that 

\ceq{\hfill\sum_{i<\lambda}\mu(\X_i)}{\le}
{\mu_*\bigg(\bigcup_{i<\lambda}\X_i\bigg)}
\medrel{\le}
$\displaystyle{\mu^*\bigg(\bigcup_{i<\lambda}\X_i\bigg)}$
\medrel{\le}
$\displaystyle{\sum_{i<\lambda}\mu(\X_i)}.$

For finite sums these inequalities hold by \ssf{\#}.
By continuity they hold also for infinite sums.
This also proves the second claim of the proposition.
\end{proof}

\section{Extension of Keisler measures}

Let $\Delta$ be a fragment. 
If $\D$ is a definable subset of $\U^{|x|}$, we denote by \emph{$\<\Delta,\D\>$\/} the fragment generated by $\Delta\cup\{\D\}$.
It consists of the sets $\big(\B\cap\D\big)\cup\big(\C\cap\neg\D\big)$ for $\B,\C$ ranging over $\Delta$.

\begin{proposition}
For every $r$ such that $\mu_*(\D)\le r\le\mu^*(\D)$ there is a Keisler measure $\mu'$ on $\<\Delta,\D\>$ that extends $\mu$.
\end{proposition}

\begin{proof}
\def\ceq#1#2#3{\parbox[t]{25ex}{$\displaystyle #1$}\medrel{#2}{$\displaystyle #3$}}
Assume $r=\mu^*(\D)$ and define 

\ceq{\hfill\mu'\big((\B\cap\D)\cup(\C\cap\neg\D)\big)}{=}{\mu_*(\B\cap\D)+\cdot\mu^*(\C\cap\neg\D)}


\end{proof}




\section{Smooth measures}

\begin{definition}
Let $\mu$ be a Keisler measure on $\U^{|x|}$ over $A$.
For $\phi(x\,;z)\in L$ we say that $\mu$ is \emph{$\phi(x\,;z)$-smooth\/} if $\phi(\U\,;b)$ is measurable for all $b\in\U^{|z|}$.
We say that $\mu$ is \emph{smooth\/} if it is $\phi(x\,;z)$-smooth for all  $\phi(x\,;z)\in L$.\QED
\end{definition}

Every measurable set $\X$ can be approximated in measure by pair of sets $\F\subseteq\X\subseteq\O$ that are type-definable, respectively type-codefinable, over $A$.
When $\X$ is definable (in $\U$), compactness allows us to require that $\F$ and $\O$ are both definable over $A$.
The following proposition is a uniform version of this remark.

\begin{proposition}
For every $\phi(x\,;z)\in L(A)$ the following are equivalent
\begin{itemize}
\item[1.] $\phi(\U\,;b)$ is measurable for every $b\in\U^{|z|}$;
\item[2.] for every $\epsilon>0$ there are finitely many $A$-definable sets $\C_i,\D_i\subseteq\U^{|x|}$ for $i<n$ such that
\begin{itemize}
\item[i.] $\mu(\D_i\sm\C_i)<\epsilon$;
\item[ii.] for every $b\in\U^{|z|}$ there is an $i<n$ such that $\C_i\subseteq\phi(\U\,;b)\subseteq\D_i$.
\end{itemize}
\end{itemize}
\end{proposition}
\begin{proof}
\ssf{1}$\IMP$\ssf{2} Immediate.

\ssf{1}$\IMP$\ssf{2} As observed above, for each $b\in\U^{|z|}$ there are some $A$-definable sets $\C,\D$ such that $\C\subseteq\phi(\U\,;b)\subseteq\D$ and $\mu(\D\sm\C)<\epsilon$.
Then the proposition follows by compactness.
\end{proof}
\def\Fr{{\rm Fr}}

\begin{definition}
Let  $\phi(x\;z)\in L$ and l$\mu$ be a global Keisler measure. We say that $\mu$ is \emph{$\phi(x\;z)$-fim\/} if for every $\epsilon>0$ there is are some formulas $\theta_n(x^{(n)})$ such that 
\begin{itemize}
\item[1.] $\displaystyle\lim_{n\to\infty}\mu^{(n)}\big(\theta_n(x^{(n)})\big)$
\item[2.] $a^{(n)}\models\theta_n(x^{(n)})$ 
\item[3.] $\mu\big(\phi(x\,;b)\big)$ is within $\epsilon$ of $\/\Fr_n\big(\phi(x\,;b),a\big))$
\end{itemize}

\end{definition}


\section{Measure extension}







\begin{comment}
If $A$ is a set of parameters and $x$ a variable, we let $\mathcal L_x(A)$ denote the algebra of $A$-definable sets in the variable $x$. Equivalently, it is the Boolean algebra of formulas with parameters in $A$ and free variable $x$, quotiented by the equivalence relation $\phi(x) \sim \psi(x) \iff \monster \models \phi(x) \leftrightarrow \psi(x)$. By an abuse of notations, $\phi(x)$ will be used to denote a formula as well as its image in $\mathcal L_x(A)$.

\begin{defi}
\index{Keisler (measure)}
Let $A\subset \monster$ be a set of parameters. A \emph{Keisler measure} (or simply a \emph{measure}) $\mu$ over $A$ in the variable $x$ is a finitely additive probability measure on $\mathcal L_x(A)$. In other words it is a function $\mu : \mathcal L_x(A) \rightarrow [0,1]$ such that:

$\clist$ $\mu(x=x) = 1$;

$\clist$ $\mu(\neg \phi(x))=1-\mu(\phi(x))$;

$\clist$ $\mu(\phi(x) \wedge \psi(x))+\mu(\phi(x) \vee \psi(x))=\mu(\phi(x)) + \mu(\psi(x))$.
\end{defi}
We will sometimes write $\mu$ as $\mu_x$ or $\mu(x)$ to emphasize that $\mu$ is a measure on the variable $x$. If $A \subseteq B$ and $\mu$ is a measure over $B$, we define the restriction of $\mu$ to $A$ denoted $\mu|_A$ or $\mu\upharpoonright A$ as the restriction of $\mu$ to $\mathcal L_x(A)$. As in the case of types, we say that $\mu$ is an extension to $B$ of $\mu|_A$.

\begin{ex}\label{ex_measures}
\begin{itemize}
\item A type $p\in S_x(A)$ can (and will) be identified with the Keisler measure $\mu_p(x)$ over $A$ defined by $\mu_p(\phi(x))= 1$ if $p\vdash \phi(x)$ and $\mu_p(\phi(x))=0$ otherwise. Thus a (complete) type is  a special case of a Keisler measure. We will usually not distinguish between $p$ and $\mu_p$ and write for example $p(\phi(x;b))$ instead of $\mu_p(\phi(x;b))$.

\item Given $a_0,a_1,\ldots$ in $[0,1]$ such that $\sum a_i = 1$, and types $p_0,p_1,\ldots$ over $A$ in the same variable $x$,  we can define the average measure $\mu = \sum a_i p_i$.
\item Take $T$ to be the theory of real closed fields and let $\mathbb R$ be the standard model. Let $\mu_0$ be any Borel probability measure on $\mathbb R$. Then $\mu_0$ induces a Keisler measure over $\monster$ in one variable $x$ defined by $\mu(\phi(x))=\mu_0(\phi(\mathbb R))$.
\item Let $I=(a_i : i\in [0,1])$ be an indiscernible sequence. Let $\lambda_0$ denote the usual Lebesgue measure on the interval $[0,1]$. We can define the \emph{average measure} $\Av(I)$ as the measure $\mu$ defined by $\mu(\phi(x;b))=\lambda_0(\{i\in [0,1] : \models \phi(a_i;b)\})$. Note that NIP ensures that the set in question is Lebesgue measurable (it is in fact a finite union of intervals).
\item Let $(M_n : n<\omega)$ be a sequence of finite structures. For each $n<\omega$, let $\mu_n$ denote the normalized counting measure on $M_n$. Let $\mathcal D$ be an ultrafilter on $\omega$ and consider the ultraproduct $M=\prod_{\mathcal D} M_n$. We define a measure $\mu$ on $M$ in the following way. Let $\phi(x;b)\in L(M)$ be a formula. Let $(b_n:n<\omega)$ be a representative of $b$ in $\prod_{n<\omega} M_n$. Then set $\mu(\phi(x;b))=\lim_{\mathcal D} \mu_n(\phi(M_n;b_n))$\footnote{We recall that the limit $\lim_{\mathcal D} a_i$ of a sequence $(a_i)_{i<\omega}$ of elements of a compact set $C$ is the unique $l\in C$ such that for any neighborhood $U$ of $l$, the set $\{i<\omega : a_i\in U\}$ is in $\mathcal D$.}.

\item Here is an example with IP which does not have all the nice properties that we will establish for measure in NIP theories. Let $T$ be the theory of the random graph in the language $\{R\}$. Let $M\models T$. Define a Keisler measure $\mu$ on $M$ by $\mu\left (\bigcap_{i<n} (xRa_i)^{\eta(i)}\right ) = 2^{-n}$ for any choice of pairwise distinct $a_i$'s in $M$ and $\eta : n\rightarrow \{0,1\}$.
\end{itemize}
\end{ex}

Let $\mathfrak M_x(A)$ denote the set of Keisler measures over $A$. It is a closed subset of $[0,1]^{\mathcal L_x(A)}$, equipped with the product topology. We equip $\mathfrak M_x(A)$ with the induced topology, making it a compact Hausdorff space. The identification of a type with the measure it defines gives an identification of $S_x(A)$ as a closed subspace of $\mathfrak M_x(A)$.

\subsection*{Borel measures}

\index{measures!Borel}
Let $\mu\in \mathfrak M_x(A)$ be a Keisler measure. It assigns a measure to every clopen set of the space $S_x(A)$. We show how to extend that measure to a $\sigma$-additive Borel probability measure. First, if $O\subseteq S_x(A)$ is open, we define $\mu(O)=\sup\{\mu(D) : D\subseteq O, D\text{ clopen}\}$. Similarly, the measure of a closed set $F$ is the infimum of the measures of clopen sets which contain it. If $F\subseteq O$ are respectively closed and open, then there is a definable set between them. This implies that if $X$ is either closed or open, we have $$(Reg)\quad \sup\{\mu(F) : F\subseteq X, F \text{ closed}\} = \inf\{\mu(O):X\subseteq  O, O\text{ open}\}.$$

The next step is to show that the set of subsets $X\subseteq S_x(A)$ satisfing $(Reg)$ is closed under complement and countable union. Complement is clear. For countable union: let $X=\bigcup_{i<\omega} X_i$ and fix $\epsilon>0$. For each $i<\omega$, take $F_i \subseteq X_i \subseteq O_i$ with $\mu(O_i)-\mu(F_i)\leq \epsilon 2^{-i}$. Let $O=\bigcup_{i<\omega}O_i$. Then we can find some finite $N$ such that $\mu(O)-\mu( \bigcup_{i<N} O_i) \leq \epsilon$. Let $F=\bigcup_{i<N} F_i$. Then we have $F\subseteq X \subseteq O$ and $\mu(O)-\mu(F)\leq 3\epsilon$.

It follows that every Borel subset of $S_x(A)$ satisfies $(Reg)$. We can therefore define $\mu$ on all such sets by $\mu(X)=\sup\{\mu(F) : F\subseteq X, F \text{ closed}\} = \inf\{\mu(O):X\subseteq  O, O\text{ open}\}$. It is easy to check that this defines a $\sigma$-additive measure on $S_x(A)$. Property $(Reg)$ is refered to as \emph{regularity} of the measure $\mu$.

To a Keisler measure on $A$, we have associated a regular probability measure on $S_x(A)$. Conversely, if $\mu$ is a regular probability measure on $S_x(A)$, then it defines a Keisler measure by restriction to the clopens. Regularity ensures that $\mu$ is entirely determined by that restriction. We therefore obtain a bijection: $$\{\text{Keisler measures on }A\} \longleftrightarrow \{\text{Regular Borel prob. measures on }S_x(A)\}.$$

We will always use the same notation for the Keisler measure and for the associated Borel measure on the type space. This gives meaning to the notation $\mu(X)$, where $\mu$ is a Keisler measure on $A$ and $X$ is a Borel subset of $S_x(A)$.


%Let $(\phi_n(x):n<\omega)$ is a family of disjoint non-empty $A$-formula, that we see also as a family of clopens of $S_x(A)$. Then the union $\bigcup_{n<\omega} \phi_n(x)$ is an open set of $S_x(A)$ and by compactness, is not clopen. Therefore we can apply the Hahn-Kolmogorov theorem (see for example \cite{proba}) and conclude that $\mu$ extends uniquely to a Borel $\sigma$-additive probability measure on $S_x(A)$, that we will also denote by $\mu$.

%We show that $\mu$ is \emph{regular}, which means that for any Borel set $X\subseteq S_x(A)$, the measure $\mu(X)$ is equal to the infimum of $\mu(O)$ for $X\subseteq O$, and $O$ and open set. First note that if $(\phi_n(x):n<\omega)$ is an increasing sequence of $A$-formulas, then necessarily $\mu(\bigcup_{n<\omega} \phi_n(x)) = \sup_{n<\omega} \phi_n(x)$. So measures of open subsets of $S_x(A)$ are easily obtained, and similarly for closed sets.

\smallskip
\index{support}\index{weakly random}
We define the \emph{support} of $\mu$ as the set $S(\mu)\subseteq S_x(A)$ of types $p\in S_x(A)$ such that for any $\phi(x;b)\in L(A)$, if $p(\phi(x;b))=1$ then $\mu(\phi(x;b))>0$. The support of $\mu$ is thus a closed set of $S_x(A)$. A type in the support of $\mu$ is called \emph{weakly random} for $\mu$.

\smallskip
If $X$ is a Borel set of positive $\mu$-measure, the \emph{localization} of $\mu$ at $X$ is the measure $\mu_X(x)$ defined by $\mu_X(\phi(x))=\mu(\phi(x) \cap X)/\mu(X)$.
\subsection*{Extending measures}

The next lemma says that any partial measure in a suitable sense extends to a full Keisler measure. In the statement, $\top$ denotes the true formula $x=x$.

\begin{lemme}\label{boolpartial}
Let $\Omega \subseteq \mathcal L_x(A)$ be a set of (equivalence classes of) formulas closed under intersection, union and complement, and containing $\top$. Let $\mu_0$ be a finitely additive measure on $\Omega$ with values in $[0,1]$ such that $\mu_0(\top)=1$. Then $\mu$ extends to a Keisler measure over $A$.
\end{lemme}
\begin{proof}
By compactness in the space $[0,1]^{\mathcal L_x(A)}$, it is enough to show that given $\psi_1(x),\ldots,\psi_n(x)$ definable sets in $\mathcal L_x(A)$, there is a function $f=\langle \psi_1,\ldots,\psi_n \rangle \rightarrow [0,1]$ finitely additive and compatible with $\mu_0$. (Here $\langle X \rangle$ denotes the Boolean algebra generated by $X$.) We may assume that $\psi_1,\ldots,\psi_n$ are the atoms of the Boolean algebra $B$ that they generate.

The elements of $\Omega$ in $B$ form a sub-Boolean algebra. Let $\phi_1,\ldots,\phi_m$ be its atoms. We have say:
$$ \phi_1 = \psi_{i_1(1)} \vee \cdots \vee \psi_{i_1(l_1)} \quad \phi_2 = \psi_{i_2(1)} \vee \cdots \vee \psi_{i_2(l_2)} \quad \text{etc.}$$
No $\psi_i$ appears in more than one such expression since the $\phi_i$'s are disjoint. Then any finitely additive $f$ satisfying $f(\psi_{i_1(1)})+\cdots+f(\psi_{i_1(l_1)}) = \mu_0(\phi_1)$ etc. will do.
\end{proof}

\begin{lemme}\label{lem_extmeasure}
Let $\mu\in \mathfrak M_x(M)$ be a measure and let $\phi(x;b)\in L(\monster)$. Let $r_1 = \sup \{ \mu(\psi(x)) : \psi(x) \in L(M), \models \psi(x) \rightarrow \phi(x;b)\}$ and $r_2= \inf  \{ \mu(\psi(x)) : \psi(x) \in L(M), \models \phi(x;b)\rightarrow \psi(x)\}$.

Then for any $r_1 \leq r \leq r_2$, there is an extension $\nu\in \mathfrak M_x(\monster)$ of $\mu$ such that $\nu(\phi(x;b))=r$.
\end{lemme}
\begin{proof}
It is enough to find some $\nu_1,\nu_2$ such that $\nu_i(\phi(x;b))=r_i$ for $i\in \{1,2\}$, since we can then consider averages of $\nu_1$ and $\nu_2$. We build $\nu_2$. Let $\Omega \subseteq \mathcal L_x(\monster)$ be the Boolean algebra generated by $\mathcal L_x(M)$ and the definable set $\phi(x;b)$. By Lemma \ref{boolpartial} it is enough to define $\nu_2$ on $\Omega$.

First assume that $r_2=0$. Then for any $\theta(x)\in L(M)$, set $\nu_2(\theta(x))=\mu(\theta(x))$ and $\nu_2 (\phi(x;b))=0$. This extends uniquely to a finitely additive measure on $\Omega$. A similar argument works for $r_2=1$.

If $r_2>0$, then let $\mu'$ be the localization of $\mu$ on the closed set $C=\bigwedge \theta(x)$ where $\theta(x)$ ranges over elements of $\mathcal L_x(M)$ such that $\models \phi(x;b)\rightarrow \theta(x)$. Let $\mu''$ be the localization of $\mu$ on the complementary open set. Then we have $\mu = r_2\mu' + (1-r_2) \mu''$. By the previous paragraph, we can extend $\mu'$ and $\mu''$ respectively to $\nu'$ and $\nu''$ such that $\nu'(\phi(x;b))=1$ and $\nu''(\phi(x;b))=0$. Then set $\nu_2=r_2 \nu' + (1-r_2) \nu''$.

The construction of $\nu_1$ follows by taking $\neg \phi(x;b)$ instead of $\phi(x;b)$.
\end{proof}

When manipulating types, it is often convenient to be able to realize them as points in a structure. The same thing cannot be done so easily with measures, but there are a couple of techniques that can be used instead. One of them is smooth measures which we will describe later. Another one consists in adding the interval $[0,1]$ as a new sort and coding the measure inside the new structure via definable maps.

More precisely, let $M$ be a structure and let $\mu(x)$ be a Keisler measure over $M$. We encode this data as a structure $$\tilde M_{\mu}=(M,[0,1],<,+,f_\phi: \phi(x;y)\in L),$$ where $M$ is equipped with its full structure, $[0,1]$ is the standard unit interval endowed with the ordering and addition modulo 1. For each formula $\phi(x;y)\in L$ the function $f_{\phi}(y)$ sends $b$ to $\mu(\phi(x;b))\in [0,1]$. Let now $\tilde N=(N,[0,1]^*,\ldots)$ be an elementary extension of $\tilde M_\mu$, where $[0,1]^*$ a (possibly) non-standard interval. Let $\st: [0,1]^* \to [0,1]$ be the standard part map, that is the map sending any element of $[0,1]^*$ to the unique real number infinitesimally close to it. From $\tilde N$ we recover a Keisler measure $\mu'(x)$ over $N$ by setting $\mu'(\phi(x;b))=\st(f_\phi(b))$.

\section{Boundedness properties}


\begin{lemme}\label{lem_indmeasure}
Let $\mu\in \mathfrak M_x(M)$ be a measure and $(b_i : i<\omega)$ an indiscernible sequence in $M$. Let $\phi(x;y)$ be a formula and $r>0$ such that $\mu(\phi(x;b_i))\geq r$ for all $i<\omega$. Then the partial type $\{ \phi(x;b_i) : i<\omega \}$ is consistent.
\end{lemme}
\begin{proof}
First we show that we can assume that the sequence $(b_i)_{i<\omega}$ is $\mu$-indiscernible, by which we mean that if $i_1<\dots<i_n <\omega$ and $j_1<\dots<j_n<\omega$, then $\mu(\phi(x;b_{i_1})\wedge\dots\wedge \phi(x;b_{i_n}))=\mu(\phi(x;b_{j_1})\wedge \dots \wedge\phi(x;b_{j_n}))$. To see this, expand $M$ to $\tilde M_\mu=(M,[0,1],\ldots)$ as described after Lemma \ref{lem_extmeasure}. Then in an elementary extension of $\tilde M_\mu$, we can find a sequence $(b'_i)_{i<\omega}$, indiscernible over $\tilde M_\mu$ in the extended language and realizing the EM-type of $(b_i)_{i<\omega}$. As $(b_i)_{i<\omega}$ is $L$-indiscernible, the two sequences have the same $L$-type. Therefore we may replace $(b_i)_{i<\omega}$ by $(b'_i)_{i<\omega}$, which has the desired property.

Now assume for a contradiction that $\{ \phi(x;b_i) : i<\omega \}$ is inconsistent. Then there is some $N$ such that $\mu(\phi(x;b_0)\wedge \dots \wedge \phi(x;b_{N-1}))=0$. Take a minimal such $N$ and let $N'=N-1$. For any integer $m$, let $\psi_m(x)=\phi(x;b_{mN'})\wedge \dots \wedge \phi(x;b_{mN'+N'-1})$. By minimality of $N$ and indiscernibility of the sequence, there is some $t>0$ such that we have $\mu(\psi_m(x))=t$ for all $m$. Also by the property of $N$, we have $\mu(\psi_m(x) \wedge \psi_{m'}(x))=0$ for $m\neq m'$.

But then we have $\mu(\psi_0(x) \vee \dots \vee \psi_{m-1}(x))=mt$, for all $m$. This contradicts the fact that the measure of the whole space is 1.
\end{proof}

Until now we have not used the assumption that $T$ is NIP. It comes into play in the following lemma which implies that measures are bounded objects.

\begin{lemme}\label{lem_bddtriangle}
Let $\mu\in \mathfrak M_x(M)$. We cannot find a sequence $(b_i : i<\omega)$ of tuples of $M$, a formula $\phi(x;y)$ and $\epsilon >0$ such that $\mu(\phi(x;b_i) \triangle \phi(x;b_j))>\epsilon$ for all $i,j <\omega$, $i\neq j$.
\end{lemme}
\begin{proof}
Assume otherwise. Then by Ramsey and compactness, we may assume that the sequence $(b_i : i<\omega)$ is indiscernible. By NIP, the partial type $\{ \phi(x;b_{2k})\triangle \phi(x;b_{2k+1}) : k<\omega\}$ is inconsistent. This contradicts the previous lemma.
\end{proof}

Let $\mu \in \mathfrak M_x(\monster)$ be a global measure. If $\phi(x)$ and $\psi(x)$ are two $M$-definable sets, set $\phi(x) \sim_{\mu} \psi(x)$ if $\mu(\phi(x)\triangle \psi(x))=0$. Then the previous lemma implies that the set of $\sim_\mu$ -equivalence classes has small cardinality (bounded by some $\kappa$ depending only on $|T|$). In particular, the support $S(\mu)$ of $\mu$ has small cardinality.

\section{Smooth measures}

We now define and investigate the notion of a smooth measure, which can be considered as an analog of a realized type.

\begin{defi}
\index{smooth}
Let $\mu\in \mathfrak M_x(M)$. We say that $\mu$ is \emph{smooth} if for every $N \supseteq M$, $\mu$ has a unique extension to an element of $\mathfrak M_x(N)$.

More generally, if $\mu \in \mathfrak M_x(N)$ and $M\subseteq N$, we say that $\mu$ is smooth over $M$ if $\mu|_M$ is smooth.
\end{defi}


\begin{lemme}\label{lem_smooth}
 Let $\mu \in \mathfrak M_{x}(M)$ be a smooth measure. Let $\phi(x,y)\in L$ and $\epsilon> 0$. Then there are formulas
$\theta_{i}^{1}(x), \theta_{i}^{2}(x)$ for $i=1,\dots,n$ and $\psi_{i}(y)$ for $i=1,\dots,n$, all over $M$ such that:

(i) the formulas $\psi_{i}(y)$ partition $y$-space;

(ii) for all $i$, if $\models \psi_{i}(b)$, then $\models \theta_{i}^{1}(x) \rightarrow \phi(x,b)\rightarrow \theta_{i}^{2}(x)$;

(iii) for each $i$, $\mu(\theta_{i}^{2}(x)) - \mu(\theta_{i}^{1}(x)) < \epsilon$.
\end{lemme}
\begin{proof} Let $b\in \monster$. Then by smoothness of $\mu$ and Lemma \ref{lem_extmeasure}, 
there are formulas $\theta_i ^{1}(x)$, $\theta_i ^{2}(x)$ over $M$, such that
\newline
(*)
$\models \theta_i ^{1}(x) \rightarrow \phi(x,b) \rightarrow \theta_ i^{2}(x) $, and $\mu(\theta_i^{2}(x))- \mu(\theta_i ^{1}(x)) < \epsilon$.
\newline
By compactness, there are
finitely many such pairs, say, $(\theta_{i}^{1}(x), \theta_{i}^{2}(x))$ such that for every $b$, one of these pairs satisfies (*).
It is then easy to find the $\psi_{i}(y)$'s.
\end{proof}

Note that conversely, the existence of formulas $\theta_i^t(x)$ and $\psi_i(y)$ with the above properties implies that the measure $\mu$ is smooth.

\begin{prop}\label{prop_smoothexist}
Let $\mu \in \mathfrak M_x(M)$ be any measure. Then there is $M\prec N$ and an extension $\mu'\in \mathfrak M_x(N)$ of $\mu$ which is smooth.
\end{prop}
\begin{proof}
Assume not. We build an increasing sequence $(M_\alpha : \alpha<|T|^+)$ of models and an associated increasing sequence of measures $\mu_\alpha \in \mathfrak M_x(M_\alpha)$. Set $M_0=M$. At limit stages, take the union. Assume $M_\alpha, \mu_\alpha$ are defined. By hypothesis, $\mu_\alpha$ is not smooth,
%Hence for some formula $\phi(x;b_\alpha)$, we have $r_2-r_1=:\epsilon_\alpha>0$, where $r_1=\sup\{\mu_{\alpha}(\psi(x)): \psi(x)\in L(M_\alpha), \models \psi(x)\rightarrow \phi(x;b_\alpha)\}$ and $r_2=\inf\{\mu_{\alpha}(\psi(x)) : \psi(x)\in L(M_\alpha), \models \phi(x;b_\alpha)\rightarrow \psi(x)\}$. By Lemma \ref{lem_extmeasure} there is an extension $\mu_{\alpha+1}$ of $\mu_\alpha$ to some model $M_{\alpha+1}$ containing $b_\alpha$ such that $\mu_{\alpha+1}(\phi_\alpha(x;b_\alpha))=r_1+\epsilon_\alpha/2$. Then for any $\theta(x)\in L(M_\alpha)$, $\mu_{\alpha+1}(\theta(x)\triangle \phi_\alpha(x;b_\alpha))\geq \epsilon_\alpha/2$.
%
%
so we can find some $\phi_\alpha(x;b_\alpha)\in L(\monster)$ and $\mu^1, \mu^2$ two extensions of $\mu_\alpha$ such that $\mu^2(\phi_\alpha(x;b_\alpha))-\mu^1(\phi_\alpha(x;b_\alpha)) =4 \epsilon_\alpha >0$. Then take $M_{\alpha+1}$ an extension of $M_\alpha$ containing $b_\alpha$ and set $\mu_{\alpha+1}=1/2 (\mu^1 + \mu^2)$. Observe that for any $\theta(x)\in M_\alpha$, either the $\mu^1$ or the $\mu^2$ measure of $\theta(x)\triangle \phi_\alpha(x;b_\alpha)$ is $\geq 2\epsilon_\alpha$, hence $\mu_{\alpha+1}(\theta(x)\triangle \phi_\alpha(x;b_\alpha))\geq \epsilon_\alpha$.

Having done the construction, we may assume that $\phi_\alpha = \phi$, and $\epsilon_\alpha= \epsilon >0$ are both constant. Let $\mu'$ be the measure $\bigcup_{\alpha<|T|^+} \mu_\alpha$. Then we have $$\mu'(\phi(x;b_\alpha) \triangle \phi(x;b_\beta))\geq \epsilon$$ for any $\alpha<\beta$. This contradicts Lemma \ref{lem_bddtriangle}.
\end{proof}

We now show that smooth measures in NIP theories can be approximated by averages of points.

If $\mu_0,\dots,\mu_{n-1}$ are measures and $\phi(x;b)$ is a formula, then the notation $\Av(\mu_0,\dots,\mu_{n-1};\phi(x;b))$ 
stands for $ \frac 1 n \sum_{k<n} \mu_k ( \phi(x;b))$.

Similarly, if $a_0,\dots,a_{n-1}$ are tuples, then $\Av(a_0,\dots,a_{n-1};\phi(x;b))$ stands for $ \frac 1 n |\{ i : \models \phi(a_i;b) \}|$. We extend this notation naturally to the case where $\phi(x;b)$ is replaced by an arbitrary Borel subset of some type space.

\begin{prop}\label{prop_smoothappr}
Let $\mu(x)$ be a global measure, smooth over $M$. Let $X$ be a Borel subset of $S_x(M)$, and $\phi(x;y)$ a formula. Fix $\epsilon>0$. Then there are $a_0,\dots,a_{n-1}\in\monster$ such that for any $b\in \monster$, $$\left | \mu(X \cap \phi(x;b)) - \Av(a_0,\dots,a_{n-1};X\cap \phi(x;b)) \right | \leq \epsilon.$$
\end{prop}
\begin{proof}
Fix formulas $\psi_i(y)$, $\theta_i^0(x)$ and $\theta_i^1(x)$, $i< m$ as given by Lemma \ref{lem_smooth}. Consider $\mu$ as a probability measure on the space $S_x(M)$. We can find types $(p_i : i<n)$ in $S_x(M)$ such that if $\lambda = \frac 1 n \sum_{i<n} p_i$ then for all $i<m$, $\lambda(\theta_i^0(x)\cap X)$ and $\lambda(\theta_i^1(x)\cap X)$ are within $\epsilon$ of $\mu(\theta_i^0(x)\cap X)$ and $\mu(\theta_i^1(x)\cap X)$ respectively. (For example, pick types $p_i$ at random according to $\mu$ and apply the weak law of large numbers.)

Now take points $(a_i : i<n)$ in $\monster$ such that $a_i \models p_i$. Set $\lambda' = \frac 1 n \sum_{i<n} \tp(a_i/\monster)$. Let $b\in \monster$ and let $i<n$ be such that $\models \psi_i(b)$. Then we have $\models \theta^0_i(x) \rightarrow \phi(x;b) \rightarrow \theta^1_i(x)$ and $\mu(\theta^1_i(x))-\mu(\theta^0_i(x))\leq \epsilon$. This implies that $\lambda(\theta^1_i(x)\cap X)-\lambda(\theta^0_i(x)\cap X) \leq 3\epsilon$ and therefore $\lambda'(\phi(x;b) \cap X)$ is within $3\epsilon$ of $\lambda(\phi(x;b) \cap X)$ and thus within $4\epsilon$ of $\mu(\phi(x;b) \cap X)$.
\end{proof}

We conclude that we can approximate any measure by averages of types.

\begin{prop}\label{prop_apprmeasure}
Let $\mu\in \mathfrak M_x(A)$ be any Keisler measure; let $\phi(x;y)\in L$ be a formula and fix $X_1,\dots,X_m \subseteq S_x(A)$ Borel subsets. Let $\epsilon >0$. Then there are types $p_0,\dots,p_{n-1} \in S_x(A)$ such that, for every $b\in A$ and every $k\leq m$:
$$\left | \mu(\phi(x;b)\cap X_k) - \Av(p_0,\dots,p_{n-1};\phi(x;b)\cap X_k) \right |\leq \epsilon.$$

Furthermore, one can impose that $p_k \in S(\mu)$ for all $k$.
\end{prop}
\begin{proof}
Let $\nu_x$ be a smooth extension of $\mu$ over a model $N \supset A$. The previous proposition works equally well with finitely many Borel subsets instead of one and gives points $a_0,\dots,a_{n-1} \in N$ such that $$\left | \nu(\phi(x;b)\cap X_k) - \Av(a_0,\dots,a_{n-1};\phi(x;b)\cap X_k) \right | \leq \epsilon$$ for all $b\in N$ and $k\leq m$. Then set $p_i = \tp(a_i/A)$.

To prove the ``furthermore" part, apply the theorem with $S(\mu)$ as one of the Borel sets and assume without loss that one instance of $\phi(x;b)$ is equivalent to $x=x$. Then at most an $\epsilon$ fraction of the $p_i$'s are not in $S(\mu)$. We can safely remove them: this adds at most an extra error of $\epsilon$ to the approximation.
\end{proof}

\begin{exo}
Give another proof of Proposition \ref{prop_apprmeasure} using the VC-theorem instead of smooth measures. Deduce that the number $n$ of types can be chosen so as to depend only on VC-dim$(\phi(x;y))$ and $\epsilon$ (see \cite[Lemma 4.8]{NIP2}).
\end{exo}

%Fix a model $M$. For any formula $\phi(x;b)\in L(\monster)$, let $\partial_{M} \phi(x;b)\subseteq S_x(M)$ be the set of types $p(x)$ such that both $p(x)\wedge \phi(x;b)$ and $p(x)\wedge \neg \phi(x;b)$ are consistent. It is a closed set of $S_x(M)$.
%
%\begin{prop}
%Let $\mu(x)$ be a measure over a model $M$. Then $\mu$ is smooth if and only if $\mu(\partial_{M}\phi(x;b))=0$ for every $\phi(x;b)\in L(\monster)$.
%\end{prop}
%\begin{proof}
%Assume that $\mu(\partial_M \phi(x;b))>0$ for some $\phi(x;b)$.
%\end{proof}

\section{Invariant measures}

We now extend a number of definitions from types to measures.

\begin{defi}
\index{invariant!measure}
Let $\mu \in \mathfrak M_x(\monster)$ be a measure, and let $A\subset \monster$. We say that $\mu$ is $A$-invariant if for every $b\equiv_A b'$, and $\phi(x;y)\in L$, we have $\mu(\phi(x;b))=\mu(\phi(x;b'))$.
\end{defi}

We define $\Lstp(A)$-invariant measures similarly.

\begin{defi}
\index{forking!for measures}\index{dividing!for measures}
Let $A\subseteq B$ and $\mu\in \mathfrak M_x(B)$. Then $\mu$ \emph{does not fork} (resp. divide) over $A$ if $\mu(\phi(x;b))=0$ for every formula $\phi(x;b)\in L(B)$ which forks (resp. divides) over $A$.
\end{defi}

\begin{prop}\label{prop_notforkmeasures}
Let $\mu\in \mathfrak M_x(\monster)$ and $A\subset \monster$. Then $\mu$ does not fork over $A$ if and only if it is $\Lstp(A)$-invariant.
\end{prop}
\begin{proof}
Let $\mu \in \mathfrak M_x(\monster)$ be $\Lstp(A)$-invariant. As it is a global measure, it is enough to show that it does not divide over $A$. Let $\phi(x;b)\in L(M)$ be such that $\mu(\phi(x;b))>0$. Let $(b_i : i<\omega)$ be an $A$-indiscernible sequence with $b_0=b$. Then by assumption $\mu(\phi(x;b_i))=\mu(\phi(x;b))$ for all $i<\omega$. By Lemma \ref{lem_indmeasure}, the partial type $\{\phi(x;b_i) : i<\omega\}$ is consistent. It follows that $\phi(x;b)$ does not divide over $A$.

Conversely, assume that $\mu$ does not fork over $A$. Then if $\Lstp(b/A)=\Lstp(b'/A)$, the formula $\phi(x;b)\triangle \phi(x;b')$ forks over $A$. We conclude that $\mu(\phi(x;b)\triangle \phi(x;b'))=0$.
\end{proof}

\begin{defi}
\index{finitely satisfiable!measure}\index{definable!measure}\index{Borel-definable}
Let $M \models T$ and let $\mu \in \mathfrak M_x(\monster)$:
\begin{itemize}
\item $\mu$ is \emph{finitely satisfiable} in $M$ if for every $\phi(x;b) \in L(\monster)$ such that $\mu(\phi(x;b))>0$, there is $a\in M$ such that $\monster \models \phi(a;b)$.
\item $\mu$ is \emph{definable} over $M$ if it is $M$-invariant and for every $\phi(x;y) \in L$, and $r\in [0,1]$, the set $\{ q\in S_y(M): \mu(\phi(x;b))< r$ for any $b\in \monster, b\models q \}$ is an open subset of $S_y(M)$.
\item $\mu$ is \emph{Borel-definable} over $M$ if it is $M$-invariant and the above set is a Borel set of $S_y(M)$.
\end{itemize}
\end{defi}

If $\mu$ is finitely satisfiable in $M$ then it is $M$-invariant. Also $\mu$ is definable over $M$ if and only if it is $M$-invariant and for any open (resp. closed) subset $X$ of $[0,1]$, the set $\{ q\in S_y(M): \mu(\phi(x;b))\in X$ for any $b\in \monster, b\models q \}$ is an open (resp. closed) subset of $S_y(M)$.

A measure $\mu$ is $M$-invariant (resp. finitely satisfiable in $M$), if and only if all types in $S(\mu)$ are $M$-invariant (resp. finitely satisfiable in $M$). This uses the fact that if $\mu$ is $M$-invariant, then $\mu(\phi(x;b)\triangle \phi(x;b'))=0$ whenever $b\equiv_M b'$, which follows from Proposition \ref{prop_notforkmeasures}.
On the other hand, definability of measures is a new phenomenon that does not translate to types.

\begin{lemme}
Let $\mu \in \mathfrak M_x(\monster)$ and $M\prec \monster$.

(i) If $\mu$ is smooth over $M$, then $\mu$ is finitely satisfiable and definable in $M$.

(ii) If $\mu$ is $M$-invariant and smooth, then it is smooth over $M$.
\end{lemme}
\begin{proof}
(i): Let $\phi(x;b)$ be such that $\mu(\phi(x;b))>0$. Then by Lemma \ref{lem_smooth}, taking $\epsilon$ small enough, there is a formula $\theta^0(x)\in L(M)$ such that $\models \theta^0(x) \rightarrow \phi(x;b)$ and $\mu(\theta^0(x))>0$. In particular, $\theta^0(x)$ is consistent, and therefore has a point in $M$. It follows that also $\phi(x;b)$ has a point in $M$.

Let $r\in [0,1]$ be such that $\mu(\phi(x;b)) < r$. Take $\epsilon>0$ so that $\mu(\phi(x;b))<r-2\epsilon$. Lemma \ref{lem_smooth} applied to $\phi(x;y)$ and $\epsilon$ gives us formulas $\psi_i^\epsilon(y)\in L(M)$, $i<n$ which partition of $y$-space. One of them, say $\psi_i^\epsilon(y)$, is satisfied by $b$. Then for any $b'$, $\psi_i^\epsilon(b')$ implies that $\mu(\phi(x;b'))\leq r-\epsilon< r$. This proves definability of $\mu$.

\smallskip

(ii): Assume that $\mu$ is $M$-invariant and smooth. Let $N\succ M$ be $|M|^+$ saturated such that $\mu$ is smooth over $N$. Fix a formula $\phi(x;y)$ and $r>0$. Then the set $\{q\in S_y(N) : \mu(\phi(x;b))\leq r $ for any $b\models q\}$ is an closed set of $S_y(N)$. Its projection to $S_y(M)$ is also closed, hence $\mu$ is definable over $M$.

Take now some $\epsilon >0$. Lemma \ref{lem_smooth} applied to $\phi(x;y)$ and $\epsilon$ gives us formulas $\psi_i(y;d)$, $\theta_i^0(x;d)$ and $\theta_i^1(x;d)$, $i<n$, where we have made the parameters explicit. By definability of $\mu$, there is an $M$-formula $\zeta(z)$ satisfied by $d$ and such that for any $d'\in \zeta(\monster)$, the formulas $\psi_i(y;d')$, $\theta_i^0(x;d')$ and $\theta_i^1(x;d')$, $i<n$ also satisfy the conclusion of the lemma. Then we can find such a $d'$ in $M$. It follows that $\mu$ is smooth over $M$.
\end{proof}

\begin{lemme}\label{lem_boreldef}
Let $p\in S(\monster)$ be some $M$-invariant type, then $p$ is Borel-definable over $M$.
\end{lemme}
\begin{proof}
Let $\phi(x;y)$ be a formula. We need to see that the set $\{b\in \monster : p\vdash \phi(x;b)\}$ is Borel over $M$. Let $b\in \monster$. Take a maximal $N$ such that there is a sequence $(a_i : i<N)\models p^{(N)}|_M$ with $$(*)\quad \neg (\phi(a_i;b)\leftrightarrow \phi(a_{i+1};b))\text{ for }i<N-1.$$ By the discussion on eventual types in Section \ref{sec_evtypes}, the formula $\phi(a_{N-1};b)$ holds if and only if $p\vdash \phi(x;b)$.

Let $A_N(b)$ say that there are $(a_i:i<N)$ such that $(*)$ holds and $\models \phi(a_{N-1};b)$. Similarly define $B_N(b)$ to hold if there is $(a_i:i<N)$ such that $(*)$ holds and $\models \neg\phi(a_{N-1};b)$. Then both $A_N(x)$ and $B_N(x)$ are type-definable sets. We have $p\vdash \phi(x;b)$ if and only if there is some integer $N\leq \alt(\phi)+1$ such that $A_N(x)\wedge \neg B_{N+1}(x)$ holds.

This shows that $\{b\in \monster : p\vdash \phi(x;b)\}$ is a finite Boolean combination of type-definable (over $M$) sets.
\end{proof}

\begin{prop}\label{prop_boreldef}
\index{Borel-definable}
Let $\mu\in \mathfrak M(\monster)$ be $M$-invariant, then $\mu$ is Borel-definable over $M$.
\end{prop}
\begin{proof}
This follows easily from the previous lemma and Proposition \ref{prop_apprmeasure} which says that a measure can be approximated by types in its support.
\end{proof}

We now have all we need to define the product of an invariant measure over another, similarly as we did for types. So let $\mu(x)\in \mathfrak M(\monster)$ be $M$-invariant and let $\lambda(y) \in \mathfrak M(\monster)$ be any measure. We define $\omega(x,y) =\mu(x) \otimes \lambda(y)$\index{$\otimes$}\index{product of measures} as a global measure in two variables by the formula:
$$\omega_{xy} (\phi(x,y;b)) = \int_{q\in S_y(N)} f(q) d\lambda_y|_N,$$
where $N$ is a small model containing $Mb$ and $f: q \mapsto \mu(\phi(x,d;b))$ for some (any) $d\models q$. Note that $f$ is a measurable function by Borel-definability of $\mu$. We need to see that this definition does not depend on the choice of $N$. So let $N_1 \subseteq N_2$ be two models containing $Mb$ and let $f_1,f_2$ be the associated functions. Let $\pi:S(N_2) \to S(N_1)$ be the canonical restriction map. Then $f_2=\pi^{*}(f_1)$. For any clopen set $B\subseteq N_1$, $\lambda|_{N_1}(B)=\lambda|_{N_2}(\pi^{-1}(B))$, hence this holds also when $B$ is Borel and then by approximating $f_1$ by step functions, we deduce that the integrals of $f_1$ and $f_2$ are the same.

Note that if $\mu=p$ and $\lambda=q$ are types, then we recover the usual product $p(x)\otimes q(y)$.

If $\lambda$ is itself $M$-invariant, then $\mu(x) \otimes \lambda(y)$ is $M$-invariant (since the whole construction is invariant under automorphisms fixing $M$).

\begin{exo}
If both $\mu$ and $\lambda$ are finitely satisfiable in $M$ (resp. $M$-definable), then $\mu(x) \otimes \lambda(y)$ is finitely satisfiable in $M$ (resp. $M$-definable).
\end{exo}

We now argue that the product of measures is associative. Let $\mu(x), \eta(y)$ and $\lambda(z)$ be three global measures, and assume that $\mu$ and $\eta$ are both $M$-invariant. Let $\phi(x,y,z)\in L(M)$, let $v=\left((\mu_x \otimes \eta_y)\otimes \lambda_z\right ) (\phi(x,y,z))$ and $w=\left (\mu_x \otimes (\eta_y\otimes \lambda_z)\right ) (\phi(x,y,z))$. We want to show that $v=w$.

First assume that $\mu=p$ is a type. Let $d\phi(y,z)$ denote the Borel subset of $S_{yz}(M)$ defined by $r\in d\phi \iff p\vdash \phi(x,b,c)$ where $(b,c) \models r$. Then $w=\eta_y \otimes \lambda_z (d\phi(y,z))=\int_{q\in S_z(M)} f(q) d\lambda_y$, where $f:q \to \eta(d\phi(y,c))$, for some $c\models q$. But this is exactly $v$.

Now if $\mu$ is an arbitrary invariant measure, then given a formula $\phi(x,y,z)$ and $\epsilon$, we can find a measure $\tilde \mu$, which is an average of types such that $|\mu(\phi(x,b,c))-\tilde \mu(\phi(x,b,c))|\leq \epsilon$ for all $b,c$. Let $\tilde v, \tilde w$ be the corresponding $v$ and $w$ with $\mu$ replaced by $\tilde \mu$. Then $\tilde v$ and $\tilde w$ are respectively at distance $\epsilon$ from $v$ and $w$. By the previous paragraph, $\tilde v=\tilde w$, and we conclude that $v=w$.


\medskip
If $\mu$ is an $M$-invariant global measure, then we can define by induction $\mu^{(1)}(x)= \mu(x)$ and $\mu^{(n+1)}(x_0,\dots,x_n) = \mu(x_0) \otimes \mu^{(n)} (x_1,\dots,x_n)$. And of course $\mu^{(\omega)}(x_0,x_1,\dots)$ is the union of $\mu^{(n)}(x_0,\dots,x_{n-1})$ for $n<\omega$.

This construction is different from the product measure in probability theory in that the space of $n$-types is not the $n$-fold product of the space of 1-types. Hence $\mu^{(n)}$ is not defined on the same space and carries more information than a usual product measure would. However, $\mu^{(n)}$ and the measure-theoretic $n$-fold product of $\mu$ agree on sets which are measurable for the product algebra. In other words, they agree on formulas which are Boolean combinations of formulas in one variable. This will allow us to use probability theoretic methods (mainly the law of large numbers) with $\mu^{(n)}$ playing the role of the product measure.

\smallskip
We now prove that a finitely satisfiable measure and a definable one commute (the case of types was done in Lemma \ref{lemme_fsdef}).

\begin{lemme}\label{lem_commutemt}
\index{commute}
Let $\mu(x), \lambda(y)$ be two global $M$-invariant measures. Assume that $\mu(x) \otimes p(y) = p(y) \otimes \mu(x)$ for any $p\in S_y(\monster)$ in the support of $\lambda$. Then $\mu$ and $\lambda$ commute.
\end{lemme}
\begin{proof}
Let $\phi(x,y)$ be a formula, and write as in the definition $(\mu_x \otimes \lambda_y)(\phi(x,y)) = \int_{q\in S_y(M)} f(q) d\lambda_x$. Fix $\epsilon>0$. Approximate that integral up to $\epsilon$ by a finite sum $\sum_{i=1,\dots,n} \lambda(X_i) c_i$, where $X_i = \{q\in S_y(M) : \mu(\phi(x,b))\in [r_i,t_i], \text{ for some }b\models q\}$ with $r_i,t_i \in [0,1]$. By Proposition \ref{prop_apprmeasure}, we can find types $q_1,\dots,q_m$ weakly random for $\lambda$, such that if $\tilde \lambda$ denotes the average $\frac 1 m \sum q_j$, then: 

$\clist$ $\tilde \lambda(X_i)$ is within $\epsilon$ of $\lambda(X_i)$, for all $i\leq n$;

$\clist$ $\tilde \lambda(\phi(a,y))$ is within $\epsilon$ of $\lambda(\phi(a,y))$ for all $a\in \monster$.

The first condition ensures that $\mu_x \otimes \tilde \lambda_y (\phi(x,y))$ is within $\epsilon$ of $\mu_x \otimes \lambda_y (\phi(x,y))$ and the second one ensures that $\tilde \lambda_y \otimes \mu_x(\phi(x,y))$ is within $\epsilon$ of $\lambda_y \otimes \mu_x (\phi(x,y))$. As $\mu$ commutes with $\tilde \lambda$, the result follows.
\end{proof}

\begin{prop}\label{prop_commutemm}
Let $\mu(x) \in \mathfrak M(\monster)$ be definable and $\lambda(y) \in \mathfrak M(\monster)$ be finitely satisfiable, then $\mu(x) \otimes \lambda(y) = \lambda(y) \otimes \mu(x)$.
\end{prop}
\begin{proof}
By the previous lemma, we may assume that $\lambda(y) = q(y)$ is a type. Let $M$ such that both $\mu$ and $q$ are $M$-invariant. Assume for a contradiction that there is $r\in [0,1]$ and $\epsilon>0$ such that $\mu_x \otimes q_y (\phi(x;y)) < r-2\epsilon$ whereas $q_y \otimes \mu_x (\phi(x;y)) >r+2\epsilon$. By definability of $\mu$, there is a formula $\psi(y)\in L(M)$ such that $q\vdash \psi(y)$ and for all $b\in \psi(\monster)$, we have $\mu(\phi(x;b))<r-\epsilon$. Also, by Borel-definability of $q$, there is a Borel set $X\subseteq S_x(M)$ such that $q\vdash \phi(a;y)$ if and only if $\tp(a/M)\in X$.

Pick $p_1,\dots,p_n \in S_x(M)$ such that for all $b\in M$, $\Av(p_1,\dots,p_n;\phi(x;b))$ is within $\epsilon$ of $\mu(\phi(x;b))$ and also $\Av(p_1,\dots,p_n;X)$ is within $\epsilon$ of $\mu(X)$. Realize each $p_i$ by $a_i \in \monster$. By finite satisfiability of $q$, find some $b_0\in \psi(M)$ such that $\models \phi(a_i;b_0) \iff q\vdash \phi(a_i;y)$ for all $i$.

By choice of $p_i$'s and $\psi(y)$, we have $\Av(p_1,\dots,p_n ; \phi(x;b_0)) < r$. On the other hand, by the choice of $b_0$, that quantity is equal to $\frac 1 n \sum_i q(\phi(a_i;y))= \frac 1 n \left | \{i:a_i \in X\} \right |$ which is within $\epsilon$ of $\mu(X)=q_y\otimes \mu_x (\phi(x;y))>r+2\epsilon$. Contradiction.
\end{proof}

\section{Generically stable measures}

\index{generically stable!measure}
Similarly as we did for types, we define a \emph{generically stable measure}\index{generically stable!measure} as being a global measure which is both definable and finitely satisfiable (in some small model $M$). We will prove an analog of Theorem \ref{th_genstable}. First we define a new notion.

\begin{defi}
Let $\mu(x)$ be a global $M$-invariant measure. We say that $\mu$ is \emph{fim} (frequency interpretation measure) if for any formula $\phi(x;y)\in L$ and $\epsilon>0$, there is a family $(\theta_n(x_1,\dots,x_n) : n<\omega)$ of formulas in $L(M)$ such that:

$\clist$ $\lim \mu^{(n)}(\theta_n(x_1,\dots,x_n)) = 1$;

$\clist$ for any $(a_1,\dots,a_n)\in \theta_n(\monster)$, and any $b\in \monster$, $\Av(a_1,\dots,a_n;\phi(x;b))$ is within $\epsilon$ of $\mu(\phi(x;b))$.
\end{defi}

We will see later that a measure is fim if and only if it is generically stable. Left to right is easy, but the converse requires more work.

The main result needed is an adaptation of the VC-theorem from measure theory to the context of Keisler measures.

\subsection*{VC-type results}

We call a measure $\mu(x_1,\dots,x_n)$ over $A$ \emph{symmetric} if for any permutation $\sigma$ of $\{1,\dots,n\}$ and any formula $\phi(x_1,\dots,x_n)\in L(A)$, we have $\mu(\phi(x_1,\dots,x_n)) = \mu(\phi(x_{\sigma1},\dots,x_{\sigma n}))$. This definition extends naturally to measures in infinitely many variables.

Let $\phi(x;y)$ be a formula. We fix some integer $n$ and write $\bar x=(x_1,\dots,x_n)$ and $\bar x'=(x_1',\dots,x_n')$. We define
$$f'_n(\bar x,\bar x';b)=| \Av(x_1,\dots,x_n;\phi(x;b)) - \Av(x'_1,\dots,x'_n;\phi(x;b)) |$$
and $$f_n(\bar x,\bar x')=\sup_{b\in \monster} f'_n(\bar x,\bar x';b).$$

\noindent
Note that for any $\epsilon>0$ the statement $f'_n(\bar x,\bar x;b) > \epsilon$ can be expressed as a first order formula $\theta_{n,\epsilon}(\bar x,\bar x';b)$. Similarly, the statement $f_n(\bar x,\bar x')> \epsilon$ can be expressed as $(\exists b) \theta_{n,\epsilon}(\bar x,\bar x';b)$.

\begin{lemme}\label{lem_vcm}
Fix a formula $\phi(x;y)$ and $\epsilon >0$. Let $\mu_n(x_1,..,x_n,x'_1,..,x'_n)$ be a symmetric measure over $\emptyset$. Then with notations as above we have $$\mu_n(f_n(\bar x,\bar x')>\epsilon) \leq 4 \pi_{\phi}(n) \exp \left(- \frac{n\epsilon^2}8\right ).$$
\end{lemme}
\begin{proof}

This is the analogue of Claim 2 inside the proof Theorem \ref{th_vc}. The proof is exactly the same: we just have to make the necessary translations and check that everything still makes sense.

Let $G=\{-1,1\}^n$. If $\phi(x;y)$ is a formula, we let $\ulcorner \phi(a;b) \urcorner$ be equal to 1 it the sentence $\phi(a;b)$ holds and 0 otherwise.

\smallskip
\underline{Claim}: We have $$\mu_n(f_n(\bar x,\bar x')>\epsilon) \leq \frac 1{2^{n-1}} \sum_{\sigma \in G} \mu_n\left (\left \{\bar x: \sup_b \frac 1 n \left |\sum_{i=1}^n \sigma_i \ulcorner \phi(x_i;b)\urcorner \right | >\epsilon/2\right \}\right ).$$

Call $D_\sigma$ the set of which we take the $\mu_n$-measure on the right-hand side. Note that it is $\emptyset$-definable. For $\sigma \in G$ let $\theta'_{\sigma}(\bar x,\bar x';y)$ be the formula expressing that
$$\left | \frac 1 n \sum_{i=1}^n \sigma_i \left ( \ulcorner \phi(x_i;y)\urcorner-\ulcorner\phi(x'_i;y)\urcorner\right ) \right |>\epsilon$$
and let $\theta_{\sigma}(\bar x,\bar x')=(\exists b)\theta'_{\sigma}(\bar x,\bar x';b)$. By symmetry of $\mu_n$, for any $\sigma \in G$ we have $\mu_n(\theta_{\sigma}(\bar x,\bar x'))=\mu_n(f_{n}(\bar x,\bar x')>\epsilon)$. Hence $$\mu_n\left (f_n(\bar x,\bar x')> \epsilon \right ) = \frac 1 {2^n} \sum_{\sigma \in G} \mu_n\left (\theta_{\sigma}(\bar x,\bar x')\right ).$$
The claim then follows exactly as Claim 1 in Theorem \ref{th_vc}.

\smallskip
The translation of the rest of the proof poses no difficulty either. For a fixed tuple $\bar a$ of length $n$ and $b\in \monster$, we let $A_b(\bar a)$ be the set of $\sigma\in G$ for which we have $\frac 1 n \left | \sum_i \sigma_i\ulcorner \phi(a_i;b)\urcorner \right | > \epsilon/2$. Then Chernoff's bound gives $$\frac 1 {2^n}|A_b(\bar a)| \leq 2 \exp \left ( - \frac {n\epsilon^2}{8}\right ).$$

The tuple $\bar a$ being fixed, there are only $\pi_{\phi}(n)$ possible values for the set $A_b(\bar a)$ as $b$ varies. Hence the cardinality of the set $A_*= \bigcup_{b} A_b(\bar a)$ is at most $e=2^{n+1}\pi_\phi(n) \exp \left ( - {n\epsilon^2}/{8}\right )$. This means that a tuple $\bar a$ can be in at most $e$ many sets $D_\sigma$; thus we can bound the sum   $\sum_{\sigma \in G} \mu_n (D_\sigma)$ by $2^n e$ and the lemma follows.

%As done there, we set $Y=(\mathbb Z_2)^n$ and consider its action on $\mathcal X^{2n}$ by flipping variables. We fix $\epsilon>0$ and define as in \ref{th_vc} the set $A_{*} \subseteq Y$ of bad permutations, which has measure at most $2\pi_{\phi}(2n)\exp(-n\epsilon^2/8)$, where $\pi_{\phi}$ is the shatter function of the family of sets $\{\phi(\monster;b):b\in \monster\}$. For $\sigma \notin A_*$, we have $\delta_n(\sigma(\bar x,\bar x'))\leq \epsilon/2$. For $\sigma \in A_*$ we have at any rate $\delta_n(\sigma(\bar x,\bar x'))\leq 1$. Thus for $n$ big enough, for any $\bar x,\bar x'$, $$2^{-n} \sum_{\sigma\in Y} \delta_n(\sigma(\bar x,\bar x')) \leq \epsilon.$$
%
%By the symmetry of $\mu_n$, $E_n$ equals the $\mu_n$-expectation of $\sup_{b} \delta_n(\sigma(\bar x,\bar x'))$ for any $\sigma \in Y$, hence it is also equal to the expectation of the average. Therefore for $n$ big enough, $E_n \leq \epsilon$.
\end{proof}

\begin{cor}\label{cor_vcm}
Let $\mu(x_1,\dots)$ be a Keisler measure in $\omega$ variables over some set $A$. Assume that $\mu|_{\emptyset}$ is symmetric. Let $\phi(x;y)\in L$ and fix $\epsilon >0$. Then there is $n$ such that $\mu(\exists y (f'_n(x_1,\dots,x_n,x_{n+1},\dots,x_{2n};y)> \epsilon))\leq\epsilon.$
\end{cor}
\begin{proof}
Apply Lemma \ref{lem_vcm} letting $\mu_n(x_1,\dots,x_{2n})$ be the restriction of $\mu$ to the variables $(x_1,\dots,x_{2n})$.
\end{proof}

\begin{prop}\label{prop_vcm}
Let $\mu$ be a global $M$-invariant measure and assume that $\mu^{(\omega)}(x_1,\dots)|_\emptyset$ is symmetric. Then $\mu$ is fim.\end{prop}
\begin{proof}
Fix a formula $\phi(x;y)$ and $\epsilon>0$ small enough. Write $\bar a=(a_1,\dots,a_n)$ and $\bar a'=(a'_1,\dots,a'_n)$. Define $\theta_n(\bar a,\bar a')$ to say that for all $b$, $\Av(a_1,\dots,a_n;\phi(x;b))$ is within $\epsilon/4$ of $\Av(a'_1,\dots,a'_n;\phi(x;b))$. By the previous corollary, for $n$ big enough we have $\mu^{(2n)}(\theta_n(x_1,\dots,x_{2n})) > 1-\epsilon$. In particular, there is $\bar a\in \monster$, such that we have $\mu^{(n)}(\theta_n(\bar a,x_1,\dots,x_n))>1-\epsilon$.

Now fix $b\in \monster$ and let $\zeta_n(x_1,\dots,x_n)$ say that $\Av(x_1,\dots,x_n;\phi(x;b))$ is within $\epsilon/4$ of $\mu(\phi(x;b))$. The law of large numbers (\ref{prop_weaklaw}) gives
$$\mu^{(n)}(\zeta_n(x_1,\dots,x_n)) \geq 1-4/n\epsilon^2.$$Taking $n$ large enough, we have $\mu^{(n)}(\zeta_n(x_1,\dots,x_n)) \geq 1/2$. It follows that $\mu^{(n)}(\theta_n(\bar a,x_1,\dots,x_n) \wedge \zeta_n(x_1,\dots,x_n))>0$. In particular the value $\Av(a_1,\dots,a_n;\phi(x;b))$ is within $\epsilon/2$ of $\mu(\phi(x;b))$.

To conclude observe that for any $\bar a'$ satisfying $\theta_n(\bar a,x_1,\ldots,x_n)$ and any $b$, $\Av(a'_1,\dots,a'_n;\phi(x;b))$ is within $\epsilon$ of $\mu(\phi(x;b))$. As $\mu^{(n)}(\theta_n(\bar a,x_1,\ldots,x_n))>1-\epsilon$, we have what we want.
\end{proof}

We can do the same incorporating a Borel set $X$.

\begin{prop}\label{prop_vcmB}
Let $\mu$ be a global $M$-invariant measure, and assume that $\mu^{(\omega)}(x_1,\dots)|_\emptyset$ is symmetric. Then for any Borel set $X$, formula $\phi(x;y)$ and $\epsilon>0$, there are $a_1,\dots,a_n\in \monster$ such that for all $b\in \monster$, $\Av(a_1,\dots,a_n; \phi(x;b)\cap X)$ is within $\epsilon$ of $\mu(\phi(x;b)\cap X)$.
\end{prop}
\begin{proof}
We need to go through the proof of Lemma \ref{lem_vcm} and of Proposition \ref{prop_vcm} replacing everywhere $\phi(x;b)$ by $\phi(x;b)\cap X$. Sets which were previously definable are now Borel, but this does not create any difficulty.
\end{proof}

\begin{rem}The statements above imply that if $\mu$ is $A$-invariant and $\mu^{(\omega)}$ is symmetric, then $\mu$ is equal to the average of any \emph{realization} of $\mu^{(\omega)}|_A$. More precisely, let $\lambda(x_1,\dots)$ be a global measure which extends $\mu^{(\omega)}|_A$. Let $\phi(x;b)\in L(\monster)$ be a formula and let $r=\mu(\phi(x;b))$. Fix $\epsilon>0$, then for all but finitely many values of $n$, we have $\lambda(\phi(x_n;b))\in (r-\epsilon,r+\epsilon)$.

This implies in particular that $\mu$ is determined by $\mu^{(\omega)}|_A$. (This is also true for arbitrary invariant measures, but we do not have the tools to prove it.)
\end{rem}
\begin{proof} Assume that we have $\lambda(\phi(x_n;b)) - r>\epsilon$ for infinitely many values of $n$. Without loss, this is true for all $n$. Now set $\omega = \mu^{(\omega)} (x'_1,x'_2,\dots) \otimes \lambda(x_1,x_2,\dots)$. Then $\omega|_\emptyset$ is symmetric. By Corollary \ref{cor_vcm}, for $n$ large enough, we have $\omega(f'_n(x_1,\dots,x_n,x'_1,\dots,x'_n;b)<\epsilon/4)>1/2$.

Let $\theta_n(x_1,\dots,x_n)$ say $\Av(x_1,\dots,x_n;\phi(x;b))-r>\epsilon/2$. Then by the weak law of large numbers, we have $\lambda(\theta_n(x_1,\dots,x_n)) \rightarrow 1$. Letting the formula $\theta'_n(x'_1,\dots,x'_n)$ say $|\Av(x'_1,\dots,x'_n;\phi(x;b))-r|<\epsilon/4$, then we also have $\mu^{(\omega)}(\theta'_n(x'_1,\dots,x'_n)) \rightarrow 1$. It follows that $\omega(\theta_n(x_1,\dots,x_n)\wedge \theta'_n(x'_1,\dots,x'_n))\rightarrow 1$. But this contradicts what we established in the previous paragraph.
\end{proof}

\subsection*{Properties of generically stable measures}

\index{generically stable!measure}
\begin{thm}\label{th_genstablem}
Let $\mu$ be a global $M$-invariant measure. Then the following are equivalent:

(i) $\mu$ is generically stable;

(ii) for any formula $\phi(x;y)\in L$ and $\epsilon>0$, there are $a_1,\ldots,a_n\in M$ such that for any $b\in M$, $\Av(a_1,\ldots,a_n;\phi(x;b))$ is within $\epsilon$ of $\mu(\phi(x;b))$.

(iii) $\mu$ is fim;

(iv) $\mu^{(\omega)}(x_1,\dots)|_M$ is symmetric;

(v) $\mu$ commutes with itself: $\mu(x) \otimes \mu(y) = \mu(y) \otimes \mu(x)$.
\end{thm}
\begin{proof}
(i) $\Rightarrow$ (v): Follows from Proposition \ref{prop_commutemm}.

(v) $\Rightarrow$ (iv): Clear by associativity of $\otimes$.

(iv) $\Rightarrow$ (iii) is Proposition \ref{prop_vcm}.

(iii) $\Rightarrow$ (ii): Clear.

(ii) $\Rightarrow$ (i): Easy.
\end{proof}

\begin{prop}
Let $\mu(x)$ be a generically stable measure, then for any invariant measure $\lambda(y)$, we have $\mu(x) \otimes \lambda(y) = \lambda(y) \otimes \mu(x)$.
\end{prop}
\begin{proof}
By Lemma \ref{lem_commutemt}, we may assume that $\lambda_y=q_y$ is an invariant type. Let $M$ be such that both $\mu$ and $q$ are invariant over $M$ and let $N\supset M$ be $|M|^+$-saturated. Take some $b\models q|_N$ in $\monster$ and let $\phi(x,y)\in L(M)$ be a formula. Let $X\subseteq S_x(M)$ be the set of types $p$ such that $q\vdash \phi(a,y)$ for some (any) $a\models p$. Then, by definition, we have $\mu_x \otimes q_y (\phi(x,y)) = \mu(\phi(x,b))$ and $q_y \otimes \mu_x(\phi(x,y)) = \mu(X)$.

Fix $\epsilon>0$. By Proposition \ref{prop_vcmB} there are $a_1,\dots,a_n\in N$ such that for all $b'\in \monster$, $\Av(a_1,\dots,a_n; \phi(x;b'))$ is within $\epsilon$ of $\mu(\phi(x;b'))$ and also $\Av(a_1,\dots,a_n;X)$ is within $\epsilon$ of $\mu(X)$. Then $\mu_x \otimes q_y (\phi(x,y))$ is equal to $\Av(a_1,\ldots,a_n;X)$ hence within $\epsilon$ of $\Av(a_1,\dots,a_n; \phi(x;b))$, which by definition of $X$ is within $\epsilon$ of $q_y \otimes \mu_x (\phi(x,y))$.

As this holds for all $\epsilon>0$, the result follows.
\end{proof}

\begin{prop}
Let $\mu(x)$ be generically stable and $A$-invariant. Then $\mu$ is the unique $A$-invariant extension of $\mu|_A$.
\end{prop}
\begin{proof}
Assume that $\lambda(x)$ is a global $A$-invariant extension of $\mu$. As in the proof of Proposition \ref{prop_uniqueext}, we show that $\mu^{(\omega)}|_A = \lambda^{(\omega)}|_A$. It follows that $\lambda$ is also generically stable, and by the remark after Proposition \ref{prop_vcmB}, $\mu=\lambda$.
\end{proof}

\begin{ex}\label{ex_mesuresgenstable}
Recall the examples of measures given in Example \ref{ex_measures}. We go through them again, and see that most of them are generically stable.
\begin{itemize}
\item If $p\in S_x(\monster)$ is a global type, then $p$ is generically stable as a measure if and only if it is generically stable as a type. More generally, an average $\sum_{i<\omega} a_ip_i$ of types is generically stable if and only if all the $p_i$'s are generically stable (assuming of course $a_i>0$ for all $i$).
\item Any Borel probability measure on $\mathbb R$ induces a smooth Keisler measure. See \cite{NIP3}. In fact, more generally, under mild assumptions, if $M$ is a model  equipped with a $\sigma$-algebra making externally definable sets measurable and $\mu_0$ is a $\sigma$-additive measure on $M$, then $\mu_0$ induces a generically stable Keisler measure. See \cite[Theorem 5.1]{finding} for a precise statement and a proof.
\item The average measure of an indiscernible sequence $I=(a_i : i\in [0,1])$ is always generically stable. It is easy for example to check fim. It is smooth if and only if the sequence $I$ satisfies a property called \emph{distality}. See Proposition \ref{prop_distalseqmeasures}.
\item Let $M=\prod_{\mathcal D} M_n$ be an ultraproduct of finite structures. Define $\mu_n$ as the normalized counting measure on $M_n$ and let $\mu$ be an ultralimit of the $\mu_n$'s. Then $\mu$ is generically stable. Indeed the VC-theorem implies that property (ii) in Theorem \ref{th_genstablem} holds for $\mu_n$ where the number of points in the approximation depends only on $\phi(x;y)$ and $\epsilon$. Hence this is also true at the limit.
\end{itemize}
\end{ex}

\begin{exo}
Let $\mu$ be the average of the indiscernible sequence $I=(a_i:i<\in [0,1])$. Given $n<\omega$ and a formula $\phi(x_1,\ldots,x_n)\in L(\monster)$ give an explicit formula for $\mu^{(n)}(\phi(x_1,\ldots,x_n))$ in terms of the sequence $(a_i)$. Check that $\mu^{(n)}$ is symmetric, which gives another proof that $\mu$ is generically stable.
\end{exo}

\begin{exo}
Let $\mu(x)$ be $M$-invariant and generically stable. Let $\phi(x)\in L(\monster)$ such that $\mu(\phi(x))>0$. Then there is some $\psi(x)\in L(M)$ such that $\psi(M)\subseteq \phi(M)$ and $\mu(\psi(x))>0$.
\end{exo}
\end{comment}
\end{document}
